\begin{lemma}
Fix \(n,m,p\), an event \(E\) in the two-bite sample space, and a
nonnegative constant \(C\).  Assume all two-bite sample weights are
nonnegative.  Suppose that for every red base graph \(G_R\) and every red row
map \(\rho\), the conditional probability of \(E\) given the cell
\[
\omega_1=G_R,\qquad v\mapsto \pi_R(v)=\rho(v)
\]
is at most \(C\).  Then the unconditional probability of \(E\) is at most
\(C\).
\end{lemma}
\begin{proof}
The cells indexed by pairs \((G_R,\rho)\) form a finite partition of the
two-bite sample space: every sample has exactly one red base graph and exactly
one red-coordinate map.  Let \(C_{G_R,\rho}\) denote such a cell and let
\(\mu\) be the two-bite weight.  Then
\[
\Pr(E)=\sum_{G_R,\rho}\mu(E\cap C_{G_R,\rho}).
\]

For a fixed cell, there are two cases.  If \(\mu(C_{G_R,\rho})=0\), then
\(\mu(E\cap C_{G_R,\rho})=0\) as well: the event \(E\cap C_{G_R,\rho}\) is a
subset of \(C_{G_R,\rho}\), and all summands in the finite weighted sum are
nonnegative, so a subset of a zero-mass cell also has zero mass.  This is
exactly the zero-mass convention in
\(\noderef{TwoBiteConditionalProbability}\): the conditional probability on
that cell is defined to be \(0\), and the cell contributes nothing to the
unconditional sum.

If \(\mu(C_{G_R,\rho})\ne0\), then the conditional probability is
\[
\frac{\mu(E\cap C_{G_R,\rho})}{\mu(C_{G_R,\rho})}.
\]
Nonnegativity of weights gives \(\mu(C_{G_R,\rho})>0\), so multiplying by the
cell mass preserves the inequality.
The hypothesis says this ratio is at most \(C\), hence
\[
\mu(E\cap C_{G_R,\rho})\le C\,\mu(C_{G_R,\rho}).
\]
Summing over all cells gives
\[
\Pr(E)\le C\sum_{G_R,\rho}\mu(C_{G_R,\rho}).
\]
The cells partition the whole sample space, whose total two-bite mass is at
most \(1\) by \(\noderef{TwoBiteEventProbabilityTotalMassBound}\).  Since
\(C\ge0\), multiplying the inequality
\(\sum_{G_R,\rho}\mu(C_{G_R,\rho})\le1\) by \(C\) gives
\[
C\sum_{G_R,\rho}\mu(C_{G_R,\rho})\le C.
\]
Thus \(\Pr(E)\le C\), as claimed.
\end{proof}
